\documentclass[11pt]{article}
\usepackage[T1]{fontenc}
\usepackage{lmodern}
\usepackage[margin=1.03in]{geometry}
\usepackage{microtype}
\usepackage{amsmath,amssymb,amsthm,mathtools,booktabs}
\usepackage[hidelinks]{hyperref}
\hypersetup{pdftitle={Sharp fifth and sixth Taylor coefficients in Hardy spaces}}
\newtheorem{theorem}{Theorem}
\newtheorem{proposition}[theorem]{Proposition}
\newtheorem{lemma}[theorem]{Lemma}
\newcommand{\D}{\mathbb D}
\newcommand{\T}{\mathbb T}
\newcommand{\Dif}{\mathcal D}
\newcommand{\coef}[1]{[z^{#1}]}
\allowdisplaybreaks[1]
\title{Sharp fifth and sixth Taylor coefficients\\
in Hardy spaces below \texorpdfstring{$H^1$}{H1}}
\author{}
\date{27 November 2026}
\begin{document}
\maketitle

\begin{abstract}
For the classical Hardy space on the disk, let \(C(n,p)\) denote
the sharp norm of the \(n\)th Taylor coefficient. We give a finite
formula for \(C(5,p)\) throughout \(0<p\leq1/2\) with at most seventy
explicit algebraic candidates. We determine \(C(5,1/2)\),
\(C(5,2/3)\), and \(C(6,2/3)\) by isolated real roots and classify
their extremizers. The numbers \(C(5,2/3)\) and \(C(6,2/3)\) have
degrees 16 and 20, respectively, and neither is expressible
by radicals. For every fixed order, we prove uniqueness of the
extremizer modulo rotations for sufficiently small \(p\), and give
two correction terms in the small-exponent asymptotic. The proofs use a
reflected stationarity identity, phase variations, and exact
polynomial elimination.
\end{abstract}

\section{The coefficient problem}

Write \(m\) for normalized arc measure on \(\T\). For \(p>0\), let
\[
 \|f\|_p=\sup_{0<r<1}
 \left(\int_{\T}|f(rz)|^p\,dm(z)\right)^{1/p},\qquad
 C(n,p)=\sup_{\|f\|_p\leq1}|[z^n]f| ,
\]
where the supremum is over holomorphic functions on \(\D\). Zeros
are allowed in this definition. Brevig and Saksman
\cite{BrevigSaksman} posed the sharp coefficient problem for \(0<p<1\)
and computed \(C(2,p)\) and \(C(3,2/3)\). The outer theorem of Brevig
and S{\o}vik \cite{BrevigSovik} gives a structural reduction for
every order.

The polynomial-jet approach to low coefficient orders draws on
the zero-free coefficient work in the Verifiable Mathematical
Reasoning Corpus \cite{VMRC}. The extremal problem here concerns
the entire Hardy space. Outer factorization and boundary
separation allow direct variations of complex polynomial
coefficients. In order five, independent rotations of groups of
coefficients give positivity on a full range of exponents. At
\(p=2/3\), positivity and polynomial elimination select one
stationary point in each of orders five and six. The resulting
Galois groups explain why the constants cannot have radical
formulas. The small-\(p\) argument applies in every fixed order.

\section{Boundary separation and stationary jets}





\begin{proposition}\label{prop:boundary}
Fix \(n\ge1\), \(0<p<1\), and \(s=2/p\). The norm \(C(n,p)\)
is attained. Every extremizer of quasi-norm one has the form
\[
 f=\eta Q^s,\qquad
 \deg Q=n,\quad \|Q\|_2=1,\quad Q(0)>0,\quad
 Q(z)\ne0\ \text{for }|z|\le1,\quad |\eta|=1 .
\]
The power uses the analytic logarithm of \(Q\) which is real at zero.
\end{proposition}

\begin{proof}
The Hardy pointwise estimates give a normal family on every smaller
disk. A locally uniform limit of a maximizing sequence still has
quasi-norm at most one, by integration on each smaller circle.
The coefficient converges and the maximum is positive, since \(z^n\)
is admissible. Scaling shows every extremizer has quasi-norm one.
Brevig--S{\o}vik's outer theorem \cite[Thm.~1]{BrevigSovik} applies
to the bounded linear functional \(\coef n\). For convenience, its
concavity argument is as follows. Write \(f=IO\), where \(I\) is
inner and \(O\) is outer. If \(I\) is nonconstant, set
\(u=2\Re I-c\) on \(\T\), with
\(c=\int_{\T}2\Re I\,|f|^p\,dm\). Then \(u\) is nonconstant and has
mean zero under \(|f|^p\,dm\). The analytic function
\(g=O(1+I^2-cI)\) belongs to \(H^p\), and \(g=uf\) on \(\T\).
For small real \(t>0\), strict concavity of \(x^p\) gives
\(\|f\pm tg\|_p^p<1\), since \(1\pm tu>0\).
Averaging the two coefficient values contradicts extremality.
Thus \(f\) is outer. Here \(\Re I\) cannot be constant on \(\T\)
unless \(I\) is constant, and the outer boundary modulus is positive
almost everywhere.

By rotations assume \(a_0=f(0)>0\) and \(a_n=\coef n f=C(n,p)>0\).
Vary \(f e^{t g}/\|f e^{t g}\|_p\), first with an analytic polynomial
\(g\) and then with \(ig\). Differentiating the norm is justified
by boundedness of \(g\) on \(\T\). Stationarity gives
\[
 \coef n(fg)=a_n\int_{\T}|f|^p g\,dm .
\]
With \(g=z^k\) for every \(k\ge0\), Fourier uniqueness gives
\[
 h(z)=a_n+\sum_{k=1}^n
       \bigl(a_{n-k}z^{-k}+\bar a_{n-k}z^k\bigr)
     =a_n |f(z)|^p \quad\text{on }\T\text{ a.e.}
\]
The Fej\'er--Riesz theorem factors \(h=a_n |Q|^2\) with
\(Q\) free of zeros in \(\D\), \(Q(0)>0\), \(\|Q\|_2=1\), and
degree exactly \(n\) because \(a_0>0\). The normalized outer function
with boundary modulus \(|Q|^{2/p}\) is \(Q^s\), so \(f=Q^s\).

If \(U\) has degree at most \(d\), write \(U^{*d}(z)=
z^d\overline{U(1/\bar z)}\).
Let \(R=z^n h=a_n QQ^{*n}\), let \(\Dif=z\,d/dz\), and set
\(E=R\Dif\log f=s a_n Q^{*n}\Dif Q\).
The polynomial \(R\) is self-reciprocal of degree \(2n\), and
\(R=f\bmod z^{n+1}\), meaning their Taylor coefficients agree through
order \(n\). Consequently the first \(n+1\) coefficients of \(E\)
are \(k a_k\), \(0\le k\le n\). The product rule and
\((\Dif Q)^{*n}=nQ^{*n}-\Dif(Q^{*n})\) give
\[
 E-E^{*2n}=s(\Dif R-nR).
\]
Set \(\beta=1-1/p=1-s/2\). The first \(n+1\) coefficients of
\((E+E^{*2n})/2\) are \((n-\beta(n-k))a_k\).
Reflecting the remaining coefficients shows that, on \(\T\),
\begin{equation}\label{eq:boundary-identity}
 z^{-n}\frac{E+E^{*2n}}{2}=nh-\beta D_h,
\end{equation}
where \(D_h\) is the outward radial derivative of the harmonic
extension of \(h\), the Fourier multiplier \(|k|\).

Suppose \(Q(\zeta)=0\) for some \(\zeta\in\T\). Then \(E(\zeta)=0\)
by its factored form and \(E^{*2n}(\zeta)=0\) by reflection.
After rotation to \(\zeta=1\), write \(Q(z)=(z-1)B(z)\).
Applying the multiplier to
\(h=a_n(2-z-z^{-1})|B|^2\) at 1 gives
\(D_h(1)=-2a_n\|B\|_2^2<0\), since
\(2|k|-|k+1|-|k-1|\) vanishes except at \(k=0\), where it
equals \(-2\). Equation \eqref{eq:boundary-identity} would now give
\(0=-\beta D_h(\zeta)\), which is impossible. This proves strict
boundary separation.
\end{proof}

There are two useful forms of stationarity. Put \(a=s-1>1\),
\(q_0=Q(0)\), and \(P=Q/q_0\). With \(a_n>0\) as in the proof,
the leading coefficient \(d=\coef n P\) is positive because the
constant term of \(R=a_n QQ^{*n}\) is \(q_0^s>0\).
The congruence \(R=f\bmod z^{n+1}\) gives
\begin{equation}\label{eq:jet}
 P^a\equiv P^{*n}/d\pmod{z^{n+1}},\qquad
 a_n=\frac{q_0^{a-1}}{d},\qquad
 q_0=(1+|[z]P|^2+\cdots+|\coef nP|^2)^{-1/2}.
\end{equation}
When rotating the argument without the \(a_n>0\) normalization,
\(P^a\equiv P^{*n}/\bar d\bmod z^{n+1}\) and
\(\coef n Q^s=q_0^{a-1}/\bar d\). These identities also follow
from \eqref{eq:jet} by rotation.
For polynomial coefficient optimization we instead use the complex
derivatives of \(F(q)=\coef n Q^s\) as a local function of
\(q=(q_0,\ldots,q_n)\). Strict boundary separation is an open
condition on the coefficients. All sufficiently small variations
on the complex \(H^2\) unit sphere are therefore admissible.
Allowing the local analytic branch also for complex \(q_0\) near
its positive value, homogeneity gives
\begin{equation}\label{eq:sphere-lagrange}
 \frac{\partial F}{\partial q_j}=s F\,\overline{q_j},
                  \qquad 0\le j\le n
\end{equation}
at a maximum of \(|F|\). Indeed, apply Lagrange multipliers to
\(|F|^2\) and use \(\sum q_j\,\partial F/\partial q_j=sF\).
Under rotation of the argument, \(F\) may be complex.


For later use we record the converse value identity. If \(P(0)=1\),
\(d=[z^n]P\ne0\), \(P^a\equiv P^{*n}/\bar d\pmod{z^{n+1}}\), and
\(P\) is zero-free on the closed disk, then \(Q=P/\|P\|_2\) gives an
admissible \(Q^s\) of norm one. Multiplying the jet by \(P\) gives
\([z^n]Q^s=q_0^{s-2}/\bar d\) because
\([z^n]PP^{*n}=\|P\|_2^2\).

\begin{lemma}\label{lem:positive-integer}
For every \(n\geq1\) and integer \(s\geq3\), each extremal
polynomial factor for \(p=2/s\) has all its coefficients positive
after an argument rotation and normalization of its constant term.
\end{lemma}

\begin{proof}
For any polynomial \(Q\) of \(H^2\) norm one, \(Q^s\) is
admissible. Replacing the coefficients of an extremal factor by
their moduli preserves its \(H^2\) norm and can only increase the
coefficient modulus of \(Q^s\). The resulting polynomial still
gives an extremal, so it satisfies the stationary jet. Normalize
its constant term to one and call the resulting polynomial \(P\).
Let \(S\) be the support of \(P\), and \(J\) the support of \(P^a\)
through order \(n\). The constant term of \(P\) and positivity of
the expansion imply \(S\subseteq J\). The jet gives \(J=n-S\),
so \(S=J=n-S\). The quadratic term in \(P^a\) shows that \(S\)
is closed under sums which are at most \(n\). Symmetry then
implies closure under nonnegative differences. It follows by
Euclidean division that \(S\) consists of the multiples of its
smallest positive member \(h\) up to \(n\).

If \(h>1\), \(Q^s(z)=g(z^h)\) implies \(C(n,p)\leq C(n/h,p)\)
by equality of the boundary quasi-norms of \(g(z^h)\) and \(g\).
In fact \(C(k,p)<C(n,p)\) for \(0\leq k<n\), with the natural
definition when \(k=0\). Indeed multiplication of a norm-one
extremizer in order \(k\) by \(z^{n-k}\) preserves the quasi-norm
and the coefficient value, and would give an extremizer with a
zero in the disk if equality held. This contradicts outer
factorization. Therefore \(h=1\) and all the moduli are positive.

For the original complex extremizer, equality in the triangle
inequality makes all the monomials in \([z^n]Q^s\) have the same
phase. With \(q_0>0\), the monomials proportional to
\(q_jq_{n-j}q_0^{s-2}\) and
\(q_1q_{j-1}q_{n-j}q_0^{s-3}\) show that the phase of \(q_j\)
is the sum of the phases of \(q_1\) and \(q_{j-1}\), for
\(2\leq j\leq n\). This proves the lemma.
\end{proof}



\section{The fifth coefficient for \texorpdfstring{\(0<p\leq 1/2\)}{0<p<= 1/2}}
\label{sec:fifth-half-range}

We give an exact finite formula using two explicitly displayed
polynomials. One of the variables is determined by a cubic with
exactly one positive root. The procedure has at most seventy
candidate pairs for each exponent in the stated range.

\begin{theorem}\label{thm:fifth-half-range}
Assume \(0<p\leq1/2\), and put \(s=2/p\), \(a=s-1\) and
\(\delta=a-1\). For \(\lambda>0\), define
\begin{align*}
 D&=1+\lambda+\delta t\left(\lambda+\frac{a-2}{3\delta}\right),
 \qquad V=\lambda/D,\\
 A&=\frac{\delta^2(\lambda+1)}2,
\\ B&=\frac{\delta^3}{8}
   \left(\lambda^2+\frac{2(a-2)}{\delta}\lambda
                +\frac{(a-2)(a-3)}{3\delta^2}\right),\\
 U_4&=\frac{\delta^2(\lambda+1)}4\bigl(\delta\lambda+a-2\bigr),\\
 U_5&=\frac{\delta^3(a-2)\lambda^2}{8}
          +\frac{\delta^2(a-2)(a-3)\lambda}{12}
          +\frac{\delta(a-2)(a-3)(a-4)}{120}.
\end{align*}
Let \(t=t(\lambda)\) be the unique positive root of
\begin{equation}\label{eq:fifth-F}
 F(t,\lambda)=D(D-\lambda)-\lambda^2(At^2+Bt^3)=0.
\end{equation}
Take the positive \(\lambda\) for which
\begin{equation}\label{eq:fifth-G}
 G(t(\lambda),\lambda)
   =\frac{V}{a}+\delta V^2t^2+V^3(U_4t^4+U_5t^5)-1=0.
\end{equation}
There are at most seventy such \(\lambda\). For each one, set
\[
 \begin{gathered} v=\sqrt V,\qquad x=vt,\quad b=\frac{\delta\lambda x^2}{2},
 \quad c=v\left(b+\frac{\delta x^2}{2}\right),\quad
 e=vx,\quad d=v/a,\\ 
 P(z)=1+xz+bz^2+cz^3+ez^4+dz^5 .\end{gathered}
\]
Keep \(P\) if its zeros all lie outside the closed disk and set
\(q_0=(1+x^2+b^2+c^2+e^2+d^2)^{-1/2}\). Then
\[
 C(5,p)=\max_{\text{\(P\) kept}}\frac{q_0^{a-1}}{d}.
\]
The list is nonempty. All norm-one extremizers are precisely the
unimodular scalar multiples and argument rotations of
\((q_0P)^s\) arising from the maximizing candidates.
\end{theorem}

\begin{proof}
By Proposition~\ref{prop:boundary}, a normalized extremal has the
form \(Q^s\), up to unimodular phase, with \(Q=q_0P\) of degree five,
\(q_0=Q(0)>0\), and \(P\) zero-free on the closed disk. Write
\(P=1+xz+bz^2+cz^3+ez^4+dz^5\) and
\(A_j=\binom{a}{j}\). Equation~\eqref{eq:jet}, also after an
argument rotation, is
\(P^a\equiv P^{*5}/\bar d\pmod{z^6}\).
In particular the phase of the fifth coefficient of \(Q^s\) is
\(d/|d|\) and its magnitude is \(q_0^{a-1}/|d|\).

If \(x=0\), the first jet gives \(e=0\), the fourth gives
\(A_2b^2=0\), and the second then gives \(c=0\). Up to argument
rotation \(q_0,q_5>0\) in this two-coordinate case. Its stationarity
gives \(q_0^2=a/s\), \(q_5^2=1/s\).
Perturb \(q_1=q_4=\varepsilon>0\) and scale \(q_0,q_5\) by
\(\sqrt{1-2\varepsilon^2}\).
Writing \(M=sq_0^{s-1}q_5\), the coefficient of
\(\varepsilon^2\) in the gain is
\(s(aq_0^{s-2}-M)>0\), since
\(a/(s q_0 q_5)=\sqrt{a}>1\).
These local coefficient variations are admissible by strict
boundary exclusion. We conclude \(x\ne0\), and rotate to \(x>0\).
The first two jets now give
\begin{equation}\label{eq:fifth-first-jets}
 e=axd,\qquad c=d(a\bar b+A_2x^2).
\end{equation}

We next show \(b\ne0\). If \(b=0\), the third jet and
\eqref{eq:fifth-first-jets} imply \(d=-(a-2)x/(3a)<0\).
For \(a>3\), the terms in the fifth coefficient involving
\(d,e,c\) consequently all have negative phase, and the one
remaining term, \(A_4x^5/5\) after a positive rescaling, is positive.
Varying the common phase of \(d,e,c\) increases the squared
coefficient modulus, a contradiction.
For \(a=3\), we have \(d=-x/9\), \(e=3xd\), \(c=3dx^2\).
The fourth and fifth jets give
\[
 d^2(9+18x^2)=1,\qquad
 d^2(3+18x^2+9x^4)=1.
\]
These imply \(x^4=2/3\), whereas \(d=-x/9\) and the first
equality require \(x^2+2x^4=9\).

Let \(b=\rho e^{i\theta}\), \(\rho>0\). Apart from the factor
\(s q_0^s\), the fifth coefficient is
\begin{align*}
 d+a x e+(a b+A_2x^2)c+Y(\theta),\qquad
 Y(\theta)&=L e^{2i\theta}+M_1e^{i\theta}+H,\\
 L&=A_2x\rho^2,\qquad M_1=A_3x^3\rho,\qquad H=A_4x^5/5 .
\end{align*}
The powers arise directly by expanding \(P^s\) through order
five. We have \(L,M_1>0\) and \(H\geq0\). When \(H>0\),
\[
 \frac{M_1^2}{4LH}=\frac{5(a-2)}{3(a-3)}>1 .
\]
It follows that \(|Y|^2\), viewed as a polynomial in
\(X=\cos\theta\), is strictly increasing on \([-1,1]\).
Indeed the derivative is \(8LH X+2M_1(L+H)\), which is
positive also when \(H=0\).

Let \(Z(\theta)=d+axe+(ab+A_2x^2)c\) with the other coefficients
fixed. At the current phase, \eqref{eq:fifth-first-jets} shows
that its nonzero terms share the phase of \(d\), and \(Z\ne0\).
If \(Y\) is zero there, then \(\theta=\pi\bmod 2\pi\),
by strict increase of \(|Y|^2\) in \(X\). The same strict
derivative gives \(Y'(\pi)\ne0\), by expanding in \(\theta-\pi\).
At \(\pi\), the derivative of \(|Z|^2\) is zero:
the factor multiplying \(c\) is real and its derivative is
purely imaginary, and \eqref{eq:fifth-first-jets} shows that
\(c/d\) is real. Stationarity in \(\theta\) now forces
\(\Re(Z\overline{Y'})=0\).
If \(\phi\) rotates \(d,e,c\) jointly, the squared modulus
is \(|e^{i\phi}Z(\theta)+Y(\theta)|^2\).
At this point its \(\phi\phi\) second derivative is zero and
its mixed derivative is \(2\Re(iZ\overline{Y'})\ne0\).
This Hessian cannot be negative semidefinite. Thus \(Y\ne0\).

Varying the same joint phase forces \(Y\) to align with \(Z\).
The terms \(d+axe\) and \((ab+A_2x^2)c\) can be rephased
independently to stay aligned with \(Y\) for nearby \(\theta\).
If \(ab+A_2x^2=0\) at the original phase, then \(c=0\), so
this term can simply be omitted from the variation.
It follows that the local maximum in \(\theta\) must be a local
maximum of
\[
 |d+axe|+|c|\,|ab+A_2x^2|+|Y(\theta)|.
\]
The middle term is nondecreasing in \(X\), and the last term
is strictly increasing. We obtain \(\theta=0\bmod 2\pi\).
At this phase \(Y>0\), so \(d,e,c>0\) as well.
This proves that every extremizer has positive polynomial
coefficients after rotation.

Put \(v=ad\), \(t=x/v\) and \(\lambda=2b/(\delta x^2)>0\).
The third jet, divided by \(a\delta x^2/2\), gives
\[
 \delta x\left(\lambda+\frac{a-2}{3\delta}\right)
   = (1/v-v)\lambda-v ,
\]
which is \(v^2=\lambda/D\).
The fourth jet, divided by \(ax\), gives
\[
 1/v-v
 =\frac{v\delta^2x^2}{2}(\lambda+1)
    +\frac{\delta^3x^3}{8}
          \left(\lambda^2+\frac{2(a-2)}{\delta}\lambda+
                     \frac{(a-2)(a-3)}{3\delta^2}\right).
\]
Multiplying by \(v\) and substituting \(x=vt\) gives
\eqref{eq:fifth-F}. The fifth jet is
\[
 a d+2A_2(xe+bc)+3A_3(x^2c+xb^2)+4A_4 x^3b+A_5x^5=1/d.
\]
Multiplying by \(d=v/a\) and replacing the coefficients by their
displayed expressions gives \eqref{eq:fifth-G}. These steps are
reversible for positive \(t,\lambda\).

It remains to prove the finiteness assertion. At each \(\lambda>0\),
\(F(t,\lambda)=f_0+f_1t+f_2t^2-f_3t^3\) with
\(f_0,f_1,f_3>0\). Hence Descartes' rule and the endpoint signs
give exactly one positive root. At such a root
\[
 F_t=F_t-2F/t=-2f_0/t-f_1-f_3t^2<0.
\]
The root is therefore a real analytic function \(t(\lambda)\)
on \((0,\infty)\).
For \(\lambda\geq1\), set \(w=\lambda t\).
The right side of \(D(D-\lambda)=\lambda^2(At^2+Bt^3)\)
is bounded below by \(c_a\lambda(w^2+w^3)\).
The left side is at most \(C_a\lambda(1+w)^2\).
These follow from \(A\asymp_a\lambda\),
\(B\asymp_a\lambda^2\), and the explicit formula for \(D\).
Thus \(w\) stays bounded as \(\lambda\to\infty\).
In particular \(V\to1\) and
\(G(t(\lambda),\lambda)\to 1/a-1<0\).
So \(G\) is not identically zero on the positive branch.

The polynomials \(F\) and \(D^3G\) in \((t,\lambda)\) have
total degrees at most seven and ten. No positive common zero
can belong to a common irreducible factor over \(\mathbb C\):
at such a zero \(F_t\ne0\), so \(F\) has only one local factor
through the zero, of multiplicity one, and a common factor
would make \(G\) vanish on an open segment of the analytic
positive branch. This contradicts the nonzero limit just proved.
After removing common factors, B\'ezout's bound gives at most
seventy common zeros, hence at most that many positive pairs.

Every kept polynomial is admissible after taking the indicated
analytic power. The reversible jets give
\(P^a\equiv P^{*5}/d\pmod{z^6}\), so its fifth coefficient is
\(q_0^{a-1}/d\). Conversely every extremizer belongs to this
list after the prescribed rotations. Existence of an extremizer
now gives the maximum formula and the complete equality statement.
\end{proof}


\subsection{A single-root formula at \texorpdfstring{\(p=1/2\)}{p=1/2}}

\begin{theorem}
Let \(j\) be the unique zero in \((17/20,213/250)\) of
\begin{align*}
 G(j)={}&1944j^{13}-6804j^{12}-376326j^{11}-2075139j^{10}
 -3864942j^9-4389147j^8-2254878j^7\\
 &+1682847j^6+3135456j^5+1942455j^4+439212j^3
 -2109j^2-10000j-523.
\end{align*}
Let \(r\) be the unique positive zero of the cubic
\[
 H_4(j,r)=9j^4+18j^3r+9j^3-18j^2r^2-18j^2r-9jr^3
              -39jr^2-12jr-9r^3-6r^2-r .
\]
Use the parameters
\[
 X=\frac{j}{r^2(j+1)+2rj+r/3},\qquad x=\sqrt X,\quad
 t=rx,\quad d=t/3,\quad b=jX,\quad c=t(b+X),\quad e=tx.
\]
Then
\[
 C(5,1/2)=q_0^2/d=2.9029381801080328\ldots,\qquad
 q_0=(1+x^2+b^2+c^2+e^2+d^2)^{-1/2}.
\]
The norm-one extremizers are the scalar rotations and variable
rotations of \(q_0^4(1+xz+bz^2+cz^3+ez^4+dz^5)^4\).

\end{theorem}

\begin{proof}
Apply Lemma~\ref{lem:positive-integer} with \(s=4\).
Write \(P=Q/q_0=1+xz+bz^2+cz^3+ez^4+dz^5\).
The stationary jet gives
\(P^3\equiv P^*/d\pmod{z^6}\) and \(C(5,1/2)=q_0^2/d\).
Here \(P^*\) is coefficient reversal. The coefficients of degrees
one and two give \(e=tx\), \(c=t(b+x^2)\), with \(t=3d\).
Writing \(r=t/x\), \(j=b/x^2\), \(X=x^2\), the coefficient of
degree three gives the displayed formula for \(X\).
The coefficient of degree four gives
\[
 r^2X+\bigl(2r^2(j+1)+rj^2+rj\bigr)X^2=1 .
\]
Clearing denominators gives \(H_4(j,r)=0\).
The coefficient of degree five gives
\begin{equation}\label{eq:quintic-half-jet}
 \frac{r^2X}{3}+2r^2X^2+
   \bigl(2r^2j(j+1)+r^2(j+1)+rj^2\bigr)X^3=1.
\end{equation}
For the eliminant, set \(B=3r(j+1)+6j+1\).
Since \(X=3j/(rB)\), the polynomial obtained by multiplying the
left-hand side of \eqref{eq:quintic-half-jet} minus one by
\(r^2 B^3\) is
\[
 H_5=r^3jB^2+18j^2r^2 B+
     27j^3\bigl(2rj(j+1)+r(j+1)+j^2\bigr)-r^2B^3 .
\]
Direct expansion of the Sylvester determinant gives
\[
 \operatorname{Res}_r(H_4,H_5)
          =-177147j^{11}(j+1)^5G(j).
\]
Because \(j>0\), the necessary equation is \(G(j)=0\).

The coefficient list of \(G\) has three sign changes. Direct
substitution gives the sign pairs \((-,+),(+,-),(-,+)\) at the
endpoints of the intervals
\[
 (1359/10000,\,136/1000),\qquad
 (17/20,\,213/250),\qquad
 (17,\,18).
\]
Descartes' rule and these signs give exactly one root in each of
those intervals. The cubic \(H_4\) has negative cubic and
quadratic coefficients and a positive constant coefficient for
each \(j>0\). Its linear coefficient is unrestricted.
Its coefficient list therefore has one sign change, and there is
a unique positive root \(r\).

If \(1359/10000<j<136/1000\), then \(H_4(j,43/5000)<0\).
For example, upper-bound its positive monomials by their values
at \(136/1000\) and its negative monomials by their values at
\(1359/10000\). This gives the upper bound
\(-152051295419/625000000000000\). Therefore \(r<43/5000\),
which yields
\[
 X>\frac{1359/10000}
  {(43/5000)^2(1+136/1000)+2(43/5000)(136/1000)+(43/5000)/3}
 =\frac{637031250}{24796337}>25.
\]
This contradicts \(|x|<5\), since \(x\) is the sum up to sign of
the five reciprocal roots of \(P\), all inside the open unit disk
by strict boundary exclusion.

If \(17<j<18\), then \(H_4(j,10)\ge794756>0\), by bounding its
positive monomials from below at \(17\) and negative monomials
from below at \(18\). This gives \(r>10\). Since \(X<1/r^2\), the
left-hand side of \eqref{eq:quintic-half-jet} is less than
\[
 \frac13+\frac2{100}
  +\frac{2\cdot 18\cdot19+19+18^2/10}{10000}
 =\frac{64031}{150000}<1,
\]
another contradiction.

Thus the middle root necessarily occurs at a maximizer, and it
determines all normalized coefficients. Lemma~\ref{lem:positive-integer} gives the stated equality cases.
\end{proof}


\section{The fifth coefficient at \texorpdfstring{\(p=2/3\)}{p=2/3}}

\begin{theorem}

Let \(r\) be the unique real root in \((2,\infty)\) of
\[
 G(r)=r^8+16r^7+64r^6-148r^5-1348r^4-1576r^3
              +2560r^2+4752r+1792.
\]
Define
\begin{align*}
 Y&=\frac{r^2+4r+2}{r^3(r+5)},&
 t&=\frac{(r-2)(r^2+6r+4)}{2r^2(r+5)},\\
 y&=\sqrt Y,&
 x&=t/y,&
 b&=rt,&
 c&=(2b+x^2)y,&
 d&=2t .
\end{align*}
Then
\[
 C(5,2/3)=\frac{1}{y\sqrt{1+x^2+b^2+c^2+d^2+y^2}}
           =1.7861070728218036\ldots .
\]
All unit extremizers are unimodular multiples and variable rotations
of the cube of
\((1+xz+bz^2+cz^3+dz^4+yz^5)/\sqrt{1+x^2+b^2+c^2+d^2+y^2}\).

\end{theorem}

\begin{proof}
Write an extremizer as \(Q^3\) up to phase,
where \(\deg Q=5\) and \(\|Q\|_2=1\). Lemma~\ref{lem:positive-integer} gives a positive-coefficient factor
after rotation. Write its constant-term-one factor as
\(P(z)=1+xz+bz^2+cz^3+dz^4+yz^5\). The stationary identity
\(P^2\equiv P^*/y\pmod{z^6}\), with \(P^*(z)=z^5 P(1/z)\),
gives
\[
 d=2xy,\qquad c=(2b+x^2)y,\qquad
 2c+2xb=b/y,\qquad
 2d+2xc+b^2=x/y,\qquad
 2y+2xd+2bc=1/y.
\]
With \(t=xy\), \(Y=y^2\), \(r=b/t\), the last
three identities reduce respectively to
\[
 4Y+2t(1+1/r)=1,\qquad
 (4r+r^2)Y+2t=2+2/r,\qquad
 2Y+4t^2+4r^2t^2Y+2rt^3=1.
\]
The first two give the displayed \(Y,t\).
In particular \(t>0\) forces \(r>2\). Substituting in the third
equation gives \(G(r)=0\). Its coefficients have two sign
changes, and \(G(0)>0\), \(G(2)<0\), and \(G(r)>0\) for all
sufficiently large positive \(r\). Descartes' rule of signs
then gives exactly one root above \(2\). Existence of a maximizer
and the formulas now identify its normalized factor and value.

\end{proof}



\subsection{The arithmetic of the fifth coefficient}

Let \(r>2\) be the root specified above and let \(G\) denote its
octic polynomial. The splitting field of \(G\) has Galois group
\(S_8\). For verification, write \(\overline G\) for the reduction
at the indicated prime. Modulo \(3\),
\[
 \overline G=X^8+X^7+X^6-X^5-X^4-X^3+X^2+1 .
\]
In the quotient by this polynomial,
\[
 X^{3^8}=X,\qquad
 X^{3^4}-X=-X^5-X^4+X^3+X^2 .
\]
The successive monic remainders for the gcd with the last
expression, starting with its own monic associate, are
\(X^5+X^4-X^3-X^2,\ X^3+X^2-1,\ X^2-1,\ X,\ 1\).
Thus the reduction is irreducible because its degree is \(8\)
and every proper irreducible factor of a degree-eight divisor
of \(X^{3^8}-X\) has degree dividing \(4\).
The other factorizations used are
\begin{align*}
 \bmod 17:\quad
 &\overline G=(X-7)
       (X^7+6X^6+4X^5-X^4+5X^3+6X^2+X-1),\\
 \bmod 23:\quad
 &\overline G=(X+6)(X-10)(X-5)
       (X^2-10X-9)(X^3-11X^2-2X-10).
\end{align*}
For the septic factor modulo \(17\), the residue of \(X^{17^7}\)
is \(X\). The gcd with \(X^{17}-X\) is \(1\), as certified by
the following successive monic remainders:
\begin{align*}
 &X^6-2X^5+5X^4+7X^3+7X^2-4X-3,\\
 &X^5+7X^4-5X^3+6X^2-X-3,\\
 &X^4-2X^3-8X^2-2X-6,\quad X^3+3X^2-7X,\quad X^2+X+2,\quad X+5,\quad 1 .
\end{align*}
The quadratic modulo \(23\) has discriminant \(21\), a nonsquare
because \(21^{11}=-1\) in that field. For the cubic,
\(X^{23}-X\) has residue \(7X^2-8\), and the monic remainder
sequence is \(X^2-11,\ X+11,\ 1\), so the cubic is irreducible.
The reductions are squarefree. The Galois group is therefore
transitive and contains a \(7\)-cycle and an element with cycle
structure \(1^3\,2\,3\), whose cube is a transposition.
A nontrivial block system in degree \(8\) would have block and
quotient sizes \(2,4\) or \(4,2\). The \(7\)-cycle would act trivially
on both levels, which is impossible. Thus the group is primitive.
The graph of conjugates of its transposition has connected
components forming an invariant block system, so it is connected.
Its edge transpositions generate \(S_8\).

Write \(K=\mathbb Q(r)\). The exact value satisfies
\(C(5,2/3)^2=N(r)/D(r)\), where
\[
 N(X)=16X^7(X+5)^4
\]
and
\begin{align*}
 D(X)={}&9X^{11}+164X^{10}+1036X^9+2168X^8-2240X^7-10512X^6\\
 &+1072X^5+19616X^4+55936X^3+80256X^2+48704X+10240 .
\end{align*}
Its norm from \(K\) to \(\mathbb Q\) follows from the ratio of the
two resultants with \(G\):
\[
 \operatorname{Res}(G,N)=2^{88}7^{19},\qquad
 \operatorname{Res}(G,D)=2^{52}3^8\,5\,7^{12}\cdot18447154139 .
\]
The first follows directly from \(G(0)=2^8\cdot7\) and
\(G(-5)=-7^3\). The second is the integer resultant of the displayed
polynomials, or can be verified as their Sylvester determinant.
The norm has \(7\)-adic valuation \(7\). It cannot be the norm of a
rational element, which would be an eighth power, or of a square
in \(K\). The stabilizer of a point is a maximal subgroup of \(S_8\),
so the octic field \(K\) has no proper intermediate field other
than \(\mathbb Q\). It follows that \(C(5,2/3)^2\) generates \(K\),
while \(C(5,2/3)\) has degree \(16\) over \(\mathbb Q\).
In particular its value cannot be expressed by radicals.


\subsection{Continuation near \texorpdfstring{\(p=2/3\)}{p=2/3}}

\begin{proposition}\label{prop:continue-five}
On an open interval containing \(p=2/3\), the fifth coefficient
has a unique extremal factor up to the two rotations. Its positive
constant-term-one coefficients are the unique nearby real solution
of the stationary jet \eqref{eq:jet} continued from \(p=2/3\).
The coefficients and the resulting sharp value are real analytic
in \(p\) on this interval.
\end{proposition}

\begin{proof}
In the five real coefficient variables of \(P\), the jet system
at \(p=2/3\) has an invertible Jacobian. Its first two
equations solve for \(d,c\). The change to \((Y,t,r)\) is a
local diffeomorphism because \(x,b,y>0\), and the next two
equations are linear in \(Y,t\) with determinant \(-2r(r+5)\).
After solving these, the last equation reduces to a nonzero
rational factor times \(G(r)=0\). This is a simple root,
since Descartes' rule bounds positive roots counted with
multiplicity and the two sign changes already give two roots.

For clarity, the real amplitude Hessian is invertible as well.
For \(P(z)=\sum_{j=0}^5r_jz^j\), \(r_0=1\), define
\(A_j=[z^j]P^a\), \(U=\sum r_j^2\), \(d_5=r_5\),
and \(J_i=A_{5-i}-r_i/d_5\) for \(0\le i\le5\).
Thus \(J_5=0\). The amplitude gradient of
\(\log([z^5]P^s/U^{s/2})\) is a nonzero scalar multiple of
\(J_i-r_i\sum_j r_jJ_j/U\) for \(1\le i\le5\).
These five residuals recover \(J_0,\ldots,J_4\) linearly:
the fifth recovers the common scalar, the first four then
recover \(J_1,\ldots,J_4\), and \(r_0=1\) recovers \(J_0\).
This proves amplitude nondegeneracy. In the phase variables,
use \(r_1>0\) to fix the argument rotation. At \(p=2/3\),
the Hessian of \(\log|[z^5]P^3|\) is the negative of the weighted
variance of the linear monomial phases, with positive weights
proportional to the five monomial magnitudes. Agreement of all
phases forces all phase variables to vanish, by the equality
case of the cubing calculation. The phase Hessian is therefore
negative definite. The mixed amplitude--phase block vanishes
by conjugation invariance. The real analytic implicit function
theorem now gives a unique critical branch in the normalized
slice, and conjugation keeps it real.

Finally all global maximizers for \(p\) near \(2/3\) enter this
slice near the anchor. To see this, take any sequence of their
\(H^2\)-normalized factors with \(p\to2/3\).
Their coefficients have a convergent subsequence and
\(q_0\ge 2^{-5}\), since \(P=Q/q_0\) has reciprocal roots of
modulus at most one and \(\|P\|_2\le 2^5\).
The coefficient of its analytic power through degree five is
continuous in the coefficients and exponent when \(q_0>0\).
Testing against the anchor factor raised to the variable power
shows the limiting cube attains at least the anchor value.
The cube of any such norm-one limit is admissible, so equality
holds. The uniqueness at \(p=2/3\) forces the limit after rotations
to be the anchor. Hence only the implicit critical branch
can contain a maximizer for \(p\) sufficiently close to \(2/3\).
\end{proof}



\section{The exact sixth coefficient at \texorpdfstring{\(p=2/3\)}{p=2/3}}
\label{sec:sixth-at-23}

\begin{theorem}
Let \(t\) be the unique zero in \((216738/10^6,216739/10^6)\) of
\begin{align*}
 T(t)={}&6322176t^{10}-5927040t^9-2308800t^8+3141568t^7
          +54816t^6-521640t^5\\
        &{}+32724t^4+34032t^3-2520t^2-720t+45.
\end{align*}
Define
\begin{align*}
 D(t)&=65856t^6-44352t^5-5648t^4+5872t^3-132t^2-144t+9,\\
 N(t)&=7056t^5-3784t^4-732t^3+574t^2+18t-9,\\
 b&=\frac{(1-2t)N(t)}{D(t)},\\
 u&=\frac{1}{2t^2}-2-2b-\frac{2b^2}{1-2t},\\
 x&=\sqrt u,\\
 c&=\frac{2txb}{1-2t},\\
 d&=t(2b+u),\\
 e&=2tx,\\
 P(z)&=1+xz+bz^2+cz^3+dz^4+ez^5+tz^6.
\end{align*}
Then
\[
 C(6,2/3)=\frac{1}{t\sqrt{1+u+b^2+c^2+d^2+e^2+t^2}}
          =1.9153139208926150763\ldots.
\]
All unit-norm extremizers are unimodular multiples and argument
rotations of \(P^3/\|P\|_2^3\).
\end{theorem}

\begin{proof}
For \(p=2/3\), an extremizer has the form \(Q^3\) with
\(\|Q\|_2=1\). Lemma~\ref{lem:positive-integer} lets us
take all seven coefficients of \(Q\) positive after rotation.
Set \(P=Q/Q(0)\), write \(t\) for its leading coefficient, and
write \(P^*(z)=z^6\overline{P(1/\bar z)}\).
Stationarity gives
\[
 P^2=P^*/t\pmod{z^7}.
\]
Comparing the coefficients of \(z,z^2,z^3\) gives
\(e=2tx\), \(d=t(2b+x^2)\), and \(c(1-2t)=2txb\).
In particular \(0<t<1/2\). The coefficient of \(z^5\)
gives \(1=t^2(4+4b+2u+4b^2/(1-2t))\), where \(u=x^2\).
This gives the expression for \(u\) in the statement.
The \(z^4\) and \(z^6\) equations are respectively
\[
 \begin{split}
 0 &=t\left(2t(2b+u)+\frac{4tub}{1-2t}+b^2\right)-b,\\
 0 &=t\left(2t+4tu+2bt(2b+u)+\frac{4t^2ub^2}{(1-2t)^2}\right)-1.
 \end{split}
\]
Substitute the expression for \(u\) and multiply the two
equations by \((2t-1)^2\) and \((2t-1)^3\), respectively. Call the
resulting polynomials \(G(t,b)\) and \(H(t,b)\).
Their resultant and their degree-one subresultant in \(b\) are
\[
 -8t^7(2t-1)^9 T(t),\qquad
 8t^6(2t-1)^5\bigl(D(t)b+(2t-1)N(t)\bigr).
\]
Both vanish at the common zero. Since \(0<t<1/2\), the
resultant implies \(T(t)=0\).

The following brackets use exact rational endpoints
\(I_k=(k,k+1)/10^6\). The signs of \(N,D\) are evaluated only on
the four intervals lying in \((0,1/2)\).
\[
\begin{array}{r|cc|cc}
 k & T(k/10^6)&T((k+1)/10^6)&N|_{I_k}&D|_{I_k}\\ \hline
 -537657 &+&-& &\\
 -321842 &-&+& &\\
 -286789 &+&-& &\\
 -193390 &-&+& &\\
   60164 &+&-&-&+\\
  216738 &-&+&+&+\\
  366459 &+&-&+&-\\
  429350 &-&+&+&-\\
  567655 &+&-& &\\
  636806 &-&+& &
\end{array}
\]
There is a zero in every bracket, so the degree-ten polynomial
has no other zeros and exactly one zero in each bracket.
On the four admissible brackets \(D\) is nonzero, so the linear
subresultant gives \(b=(1-2t)N/D\). Only the second bracket in
\((0,1/2)\) makes \(b>0\). Thus an extremizer must realize the
unique tuple in the statement. Existence of an extremizer then
makes this tuple admissible without an additional root-location
test on \(P\). Finally,
\([z^6]P^3=[z^6]PP^*/t=\|P\|_2^2/t\).
Dividing by \(\|P\|_2^3\) gives the asserted norm.
\end{proof}

Here is a reproducible certificate for the elimination and the
interval signs. It uses only symbolic or rational arithmetic.
\begin{verbatim}
from fractions import Fraction as F
import sympy as S

t, b = S.symbols("t b")
u = 1/(2*t**2) - 2 - 2*b - 2*b**2/(1-2*t)
g = t*(2*t*(2*b+u)+4*t*u*b/(1-2*t)+b**2)-b
h = t*(2*t+4*t*u+2*b*t*(2*b+u)
       +4*t**2*u*b**2/(1-2*t)**2)-1
G = S.cancel((2*t-1)**2*g)
H = S.cancel((2*t-1)**3*h)

def horner(c, x):
    y = 0
    for a in c:
        y = y*x+a
    return y

tc = [6322176,-5927040,-2308800,3141568,54816,-521640,
      32724,34032,-2520,-720,45]
nc = [7056,-3784,-732,574,18,-9]
dc = [65856,-44352,-5648,5872,-132,-144,9]
T, N, D = [horner(c,t) for c in (tc,nc,dc)]
sr = S.subresultants(G,H,b)
assert S.expand(sr[-1]+8*t**7*(2*t-1)**9*T)==0
assert S.expand(sr[-2]-8*t**6*(2*t-1)**5*
                (D*b+(2*t-1)*N))==0

def interval_horner(c, lo, hi):
    a = b = F(0)
    for x in c:
        v = (a*lo,a*hi,b*lo,b*hi)
        a,b = min(v)+x,max(v)+x
    return a,b

klist = [-537657,-321842,-286789,-193390,60164,216738,
         366459,429350,567655,636806]
for i,k in enumerate(klist):
    lo,hi = F(k,10**6),F(k+1,10**6)
    s = (-1)**i
    assert s*horner(tc,lo)>0 and s*horner(tc,hi)<0
    if i in (4,5,6,7):
        for c,sgn in ((nc,-1 if i==4 else 1),
                      (dc,1 if i<6 else -1)):
            assert min(sgn*y for y in
                       interval_horner(c,lo,hi))>0
\end{verbatim}


\subsection{Arithmetic of the sixth coefficient}
The constant \(C(6,2/3)\) is algebraic of degree twenty and
is not expressible by radicals.

To prove this, reduce the degree-ten polynomial \(T\) modulo
\(229\) and normalize its leading coefficient to one.
In the resulting quotient ring, \(X^{229^{10}}=X\), while
\[
 \gcd(T,X^{229^2}-X)=\gcd(T,X^{229^5}-X)=1
 \qquad\text{over }\mathbb F_{229}.
\]
These are the finite-field irreducibility criteria for degree ten.
Modulo \(73\), the monic reduction factors as
\[
 (X-12)(X^2+28X+17)
 (X^7+15X^6+10X^5-21X^4+22X^3+21X^2+24X-13).
\]
The three factors are distinct and irreducible. For the degree-two
and degree-seven factors, the same criterion verifies divisibility
by \(X^{73^d}-X\) and coprimality with \(X^{73}-X\).
The Galois group is therefore transitive and contains a
permutation of cycle type \((1,2,7)\). Its powers give both a
\(7\)-cycle and a transposition. A transitive degree-ten group
with a \(7\)-cycle admits no block system with blocks of size
two or five. Indeed, an element of order seven would have to
fix each block and then to fix each point. Thus the group is
primitive. The conjugates of its transposition are the edges of
an invariant graph, whose components form a block system. The
graph is connected, and the edge transpositions generate the
full symmetric group \(S_{10}\).

Set \(A=C(6,2/3)^2\) and put \(K=\mathbb Q(t)\).
By the theorem,
\[
 A=\frac{1}{t^2\left(1+u+b^2+
           4t^2ub^2/(1-2t)^2+t^2(2b+u)^2+4t^2u+t^2\right)}.
\]
Working modulo \(229\), the inverses in this expression exist.
The remainder of \(A\) on the basis \(1,t,\ldots,t^9\) is
\[
 52t^9+81t^8+61t^7+25t^6+55t^5+19t^4
       +22t^3-61t^2-101t+107.
\]
In that same quotient, \(A^{(229^{10}-1)/2}=-1\).
Since the denominators are invertible modulo \(229\), these
computations take place in the reduction of
\(\mathbb Z_{(229)}[t]\). The nonconstant remainder shows
\(A\notin\mathbb Q\). The natural point stabilizer \(S_9\) is
maximal in \(S_{10}\), so \(K\) has no intermediate fields and
\(\mathbb Q(A)=K\). Also \(A\) is not a square in \(K\).
If \(A=\gamma^2\), multiply \(\gamma\) by the smallest power
\(229^k\), \(k\ge0\), that makes its coefficients on
\(1,t,\ldots,t^9\) integral at \(229\). For \(k>0\), this would
produce a nonzero element of the residue field whose square is
zero, which is impossible. For \(k=0\), reduction would make
\(A\) a square in the residue field, again impossible.
Consequently \([\mathbb Q(C(6,2/3)):\mathbb Q]=20\). A radical
expression for this constant would give a radical expression
for \(A\), whose normal closure has nonsolvable group \(S_{10}\).

For completeness, the finite-field checks above can be performed
by adding this code to the preceding certificate.
\begin{verbatim}
def powmod(g, k, f):
    r = P(1, f.get_modulus())
    g = g.rem(f)
    while k:
        if k % 2:
            r = (r*g).rem(f)
        g = (g*g).rem(f)
        k //= 2
    return r

def P(expr, p):
    return S.Poly(expr,t,modulus=p)

def verify_irr(f, p, prime_divisors):
    x = P(t,p)
    n = f.degree()
    assert powmod(x,p**n,f) == x.rem(f)
    for q in prime_divisors:
        assert S.gcd(f,powmod(x,p**(n//q),f)-x).degree()==0

f = P(T,73).monic()
F2 = P(t**2+28*t+17,73)
F7 = P(t**7+15*t**6+10*t**5-21*t**4
       +22*t**3+21*t**2+24*t-13,73)
assert f == P(t-12,73)*F2*F7
verify_irr(F2,73,[2])
verify_irr(F7,73,[7])

p = 229
f = P(T,p).monic()
x = P(t,p)
verify_irr(f,p,[2,5])
one = P(1,p)
def inv(g):
    return g.rem(f).invert(f)
def md(g):
    return g.rem(f)
b = md((one-2*x)*P(N,p)*inv(P(D,p)))
u = md(inv(2*x**2)-2*one-2*b
       -2*b*b*inv(one-2*x))
d = md(x*(2*b+u))
S2 = md(one+u+b*b+4*x*x*u*b*b*inv((one-2*x)**2)
        +d*d+4*x*x*u+x*x)
A = inv(x*x*S2)
assert A==P(52*t**9+81*t**8+61*t**7+25*t**6+55*t**5
            +19*t**4+22*t**3-61*t**2-101*t+107,p)
assert powmod(A,(p**10-1)//2,f)==-one
\end{verbatim}


\section{Every fixed order for sufficiently small \texorpdfstring{\(p\)}{p}}
\label{sec:small-p}

\begin{theorem}\label{thm:small-p}
Fix an integer \(n\geq1\). For all sufficiently small \(p>0\),
the unit-norm extremizer is unique up to multiplication by a
unimodular constant and rotation of its argument. Put
\(\epsilon=p/2\), \(s=1/\epsilon\). Its normalized factor has the form
\[
 Q(z)=\frac{1+\sum_{j=1}^n\epsilon^{j/2}v_j(\epsilon)z^j}
               {\bigl(1+\sum_{j=1}^n\epsilon^jv_j(\epsilon)^2\bigr)^{1/2}} ,
 \qquad
 v_j(0)=\frac{n!}{(n-j)!n^{j/2}} .
\]
The functions \(v_j\) are real analytic at zero and positive for
small positive \(\epsilon\). They are specified uniquely near the
displayed initial values by the polynomial equations
\begin{equation}\label{eq:small-polynomials}
 B_{n-j}=v_j B_n\quad (1\le j\le n),\qquad
 B_k=[z^k]\left(1+\sum_{j=1}^n\epsilon^j v_jz^j\right)^{1/\epsilon-1}.
\end{equation}
The factor \(\epsilon^{n/2}C(n,p)\) extends real analytically to zero.
In particular,
\[
 \log\left(\frac{n!C(n,p)}{(2n/(ep))^{n/2}}\right)
  =\frac{n^2-2n+2}{8}p
   -\frac{n(n^2-3n+1)}{48}p^2+O_n(p^3).
\]
\end{theorem}

\begin{proof}
Write \(Q=\sum_{j=0}^n q_jz^j\) for an extremal factor and
\(r_j=q_j/q_0\), so \(P=1+\sum_{j=1}^n r_jz^j\).
The variation \(Q(e^t z)/\|Q(e^t z)\|_2\), for real \(t\)
near zero, is admissible by boundary separation. The logarithmic
derivative of its coefficient modulus after taking the power
\(s\) is zero. Since the unnormalized coefficient is multiplied
by \(e^{nt}\), this gives the exact identity
\[
 \sum_{j=1}^n j|q_j|^2=n/s.
\]
Thus \(q_0^2\geq1-n/s\), and all \(r_j=O_n(s^{-1/2})\).
For \(A_k=[z^k]P^{s-1}\), the binomial expansion then gives
\(A_k=O_n(s^{k/2})\). On the other hand, the admissible trial
factor \(Q_0=\sqrt{1-n/s}+\sqrt{n/s}\,z\), for \(s>2n\),
gives \(C(n,p)\gg_n s^{n/2}\) by taking the \(n\)th coefficient
of \(Q_0^s\). The stationary jet therefore gives
\[
 |A_n|=C(n,p)q_0^{2-s}\gg_n s^{n/2},\qquad
 \bar r_j=A_{n-j}/A_n=O_n(s^{-j/2}).
\]
The derivative identity now gives \(|r_1|/\sqrt\epsilon\to\sqrt n\).
Rotate the argument so \(r_1>0\), and set
\(v_j=r_j\epsilon^{-j/2}\). These variables are uniformly bounded
along extremizers as \(\epsilon\downarrow0\).
Substitution of \(\sqrt\epsilon z\) in the jet gives
\(B_{n-j}=\bar v_j B_n\).
For each \(k\le n\), the expression \(B_k\) is a polynomial
in \(\epsilon,v_1,\ldots,v_n\). Indeed a term of binomial order
\(l\) has a factor \(\epsilon^k\binom{1/\epsilon-1}{l}\)
with \(l\le k\). At \(\epsilon=0\) only \(l=k\) contributes,
so \(B_k=v_1^k/k!\). Every convergent subsequence of normalized
extremizers therefore has the same limit
\(v_j(0)=n!/((n-j)!n^{j/2})\).

Use the real part of \(B_{n-1}-v_1B_n\) and, for \(j\ge2\),
both parts of \(B_{n-j}-\bar v_jB_n\), on the slice
\(v_1\in\mathbb R\). At zero the Jacobian is block triangular.
In the first equation the derivative is the derivative of
\(v_1^{n-1}(n-v_1^2)/n!\) at \(v_1=\sqrt n\), which is nonzero.
For each \(j\ge2\), the diagonal block is multiplication by
\(-B_n(0)\) composed with conjugation. The real analytic implicit
function theorem gives a unique solution branch near the limit.
All normalized global extremizers lie on it for small positive
\(\epsilon\). The equations are invariant under conjugation,
so the branch is real and the omitted imaginary part is zero.
This proves the uniqueness and the exact polynomial specification.

To compute the coefficient along this branch, write the scaled
value as
\[
 \left(1+\sum_{j=1}^n\epsilon^jv_j^2\right)^{-1/(2\epsilon)}
 [z^n]\left(1+\sum_{j=1}^n\epsilon^jv_jz^j\right)^{1/\epsilon}.
\]
Both factors extend analytically to \(\epsilon=0\).
At zero the product is \(e^{-v_1^2/2}v_1^n/n!\), which is
stationary in \(v_1\) at \(\sqrt n\) and independent of all
other \(v_j\). Thus the derivative of the value in \(\epsilon\)
at zero is the partial derivative with the \(v_j\) fixed at
their limiting values. The coefficient of \(\epsilon\) in the
logarithm of the product is
\[
 -\frac{v_2^2}{2}+\frac{v_1^4}{4}
   +\left(v_2-\frac{v_1^2}{2}\right)\frac{n(n-1)}{v_1^2}.
\]
For \(n=1\), read this with \(v_2=0\). Substitution of \(v_1=\sqrt n\)
and \(v_2=n-1\) gives \((n^2-2n+2)/4\).
For the next term, write \(X=v_1,\ Y=v_2,\ Z=v_3\), with the
missing coordinates set to zero, and let \((n)_k\) be the falling
factorial. At fixed \(X,Y,Z\), the first three coefficients of the
logarithm of the scaled objective are \(\ell_0,\ell_1,\ell_2\), where
\begin{align*}
 A_1&=\left(Y-X^2/2\right)(n)_2/X^2,\\
 A_2&=\left(Z-XY+X^3/3\right)(n)_3/X^3
       +\left(Y-X^2/2\right)^2(n)_4/(2X^4),\\
 \ell_0&=n\log X-X^2/2-\log(n!),\\
 \ell_1&=A_1-Y^2/2+X^4/4,\\
 \ell_2&=A_2-A_1^2/2-Z^2/2+X^2Y^2/2-X^6/6.
\end{align*}
These follow by retaining terms up to order three in the logarithm
of the normalized factor before raising it to the power
\(1/\epsilon\). At the limiting coordinates,
\(\ell_{0,XX}=-2\) and \(\ell_{1,Y}=0\) whenever \(Y\) is present.
As the branch maximizes the coefficient, its derivative in \(X\)
vanishes for \(\epsilon>0\). Hence the branch has
\(X'(0)=\ell_{1,X}/2\), where
\(\ell_{1,X}=(-n^2+4n-2)/\sqrt n\) at the limit. The
coefficient of \(\epsilon^2\) along the branch is therefore
\(\ell_2+\ell_{1,X}^2/4\). Substituting the limiting coordinates
gives \(-n(n^2-3n+1)/12\). Replacing \(\epsilon\) by \(p/2\)
proves the expansion.
\end{proof}



\begin{thebibliography}{9}
\bibitem{BrevigSaksman}
O.~F.~Brevig and E.~Saksman, Coefficient estimates for \(H^p\)
spaces with \(0<p<1\), \emph{Proc. Amer. Math. Soc.} \textbf{148}
(2020), 3911--3924.
\href{https://doi.org/10.1090/proc/14995}{DOI: 10.1090/proc/14995}.

\bibitem{BrevigSovik}
O.~F.~Brevig and \O.~S{\o}vik, Extremals for linear functionals on
\(H^p\) are outer when \(0<p<1\),
\href{https://arxiv.org/abs/2610.21788}{arXiv:2610.21788}, 2026.

\bibitem{VMRC}
Verifiable Mathematical Reasoning Corpus, manuscripts [248] and
[249] on sharp zero-free Hardy coefficients and the Krzy\.z problem,
\href{https://github.com/QEEDDD/Verifiable-Mathematical-Reasoning-Corpus/tree/e73bae9fd0c319129b19aacc5c317f7abf2865ef}{snapshot e73bae9f},
2026.
\end{thebibliography}
\end{document}
